\section*{Chapitre \ref{ch:m2:how-bonds-trade} --- Comment les obligations se négocient}

\begin{solution}{exo:m2:how-bonds-trade:1}
Seulement qu'elle était supérieure à USD~5~millions, à 99.50~: le plafond masque la taille exacte des
grandes transactions sur obligations investment grade.
\end{solution}

\begin{solution}{exo:m2:how-bonds-trade:2}
$10 - 5e_n$~: 10.00 points de base pour un dealer, 7.18 pour deux, 4.85 pour quatre. Sans
fuite, plus on en interroge, mieux c'est.
\end{solution}

\begin{solution}{exo:m2:how-bonds-trade:3}
Parce que le prix publié pour chaque obligation d'un panier est une allocation d'un prix unique convenu pour
l'ensemble du panier, et non un prix négocié pour cette obligation~; les lecteurs du flux doivent savoir
ne pas le traiter comme tel.
\end{solution}

\begin{solution}{exo:m2:how-bonds-trade:4}
6, 3 et 2 dealers.
\end{solution}

\begin{solution}{exo:m2:how-bonds-trade:5}
Il achète des parts d'ETF à bon prix, les rachète contre des obligations et vend ces obligations, empochant
la décote et la réduisant. Il peut hésiter parce que les valeurs publiées des obligations sont périmées et
qu'il risque de ne pas pouvoir les vendre près de ces valeurs, parce que le rachat livre des obligations
qu'il doit alors détenir ou financer, et parce que son propre bilan est contraint en période de
tension.
\end{solution}

\begin{solution}{exo:m2:how-bonds-trade:6}
Il peut n'avoir aucun intérêt pour l'obligation, aucun inventaire ni capacité à supporter le risque, le client
peut être de ceux dont les demandes précèdent souvent les mouvements de prix, la taille peut être trop grande pour être couverte,
ou il peut s'attendre à perdre l'adjudication de toute façon et n'en retirer que de l'information.
\end{solution}

\begin{solution}{exo:m2:how-bonds-trade:7}
99.449~: avec une demi-vie de dix secondes, les cotations les plus récentes, 99.47 et 99.42, dominent et
les plus anciennes comptent beaucoup moins~; la valeur aberrante reste néanmoins écartée.
\end{solution}

\begin{solution}{exo:m2:how-bonds-trade:8}
Chaque dealer sollicité apprend l'intention du client~; les perdants peuvent prendre position avant le
gagnant ou déprécier l'obligation, et le client paie cette fuite. Le gain apporté par une
offre supplémentaire diminue à mesure que l'on ajoute des dealers, alors que le coût de la fuite croît avec chacun d'eux, de sorte que
au-delà d'un certain point, davantage de dealers font monter le coût espéré.
\end{solution}

\begin{solution}{pb:m2:how-bonds-trade:1}
\textbf{1.} 10.0 points de base, USD~25\,000.
\textbf{2.} 8.18, 7.77, 7.85 et 8.66 points de base.
\textbf{3.} Trois dealers, 7.77 points de base, soit environ USD~19\,400.
\textbf{4.} Environ USD~2\,200.
\textbf{5.} L'espérance du meilleur de $n$ cours acheteurs croît de plus en plus lentement avec $n$.
\textbf{6.} Les dealers perdants qui savent qu'un vendeur est actif vendent avant lui, cessent d'acheter, ou
déprécient l'obligation et ses couvertures, de sorte que le prix auquel le gestionnaire exécute, ou sa
transaction suivante, est moins favorable.
\textbf{7.} 6 dealers à 0.5 point de base, 2 à 1.5.
\textbf{8.} Une vente importante met plus de temps à être absorbée et fait davantage bouger le prix, de sorte que ce qu'apprennent
les perdants vaut davantage pour eux.
\textbf{9.} Interroger les dealers dont l'\omterm{def:m2:how-bonds-trade:rfq}{axe} est acheteur, qui sont les plus susceptibles de
proposer le meilleur cours acheteur, permet au gestionnaire d'en solliciter moins et de moins divulguer.
\textbf{10.} Elle ajoute des enchérisseurs qui n'ont pas l'incitation d'un dealer à prendre position avant, et qui
peuvent vouloir l'obligation~; elle peut aussi diffuser l'information à davantage de parties, selon la conception de la plateforme.
\textbf{11.} Lorsqu'on vend de nombreuses obligations à la fois, en particulier de petites ou peu liquides, on peut
demander un prix unique pour la liste et payer moins par obligation que dans de nombreuses RFQ.
\textbf{12.} Le prix et une taille affichée sous la forme «~5MM$+$~» s'il s'agit d'investment grade,
dans les quinze minutes.
\textbf{13.} Les acheteurs voient qu'un bloc important a été échangé, mais pas sa taille, de sorte qu'ils ne
savent pas combien le dealer doit encore vendre.
\textbf{14.} S'il détient des parts d'ETF ou peut les vendre à découvert~; vendre des parts d'ETF ne revient pas
à vendre l'obligation, mais cela peut servir à couvrir le risque de crédit de l'obligation jusqu'à sa vente.
\textbf{15.} Là où l'obligation devrait se négocier, avant de demander des prix, afin de pouvoir juger les cours acheteurs
qu'il reçoit.
\textbf{16.} Ils tarifient le client, la taille et leur propre inventaire~; pour un client dont les
transactions précèdent les mouvements, ils proposent un cours acheteur plus bas, et ils peuvent ne répondre que pour apprendre.
\textbf{17.} Comparer les prix après chaque RFQ avec ceux de transactions comparables réalisées
avec moins de dealers, en contrôlant la taille et les mouvements de marché.
\textbf{18.} Les obligations high yield sont moins liquides et la fuite y est plus coûteuse, de sorte que moins de
dealers peut être optimal~; mais les conditions privées peuvent être plus dispersées, ce qui favorise davantage de dealers.
\textbf{19.} Résultat nommé~: \emph{le nombre optimal de dealers} est trois, pour un coût espéré de
7.77 points de base, soit environ USD~19\,400 sur USD~25~millions.
\textbf{20.} Parce que chaque dealer sollicité apprend aussi ce que le client veut faire.
\end{solution}

\begin{solution}{iq:m2:how-bonds-trade:1}
Une action se négocie en bourse, dans des carnets d'ordres publics et continus~; une obligation d'entreprise,
l'une des nombreuses émissions de la société, se négocie de gré à gré, principalement par \omterm{def:m2:how-bonds-trade:rfq}{demandes de
cotation} (RFQ) auprès de dealers, les transactions étant déclarées à TRACE après coup. La liquidité est
concentrée dans les grandes émissions récentes~; la plupart des obligations s'échangent rarement.
\iqlookfor{OTC et RFQ, fragmentation par émission, transparence post-négociation.}
\end{solution}

\begin{solution}{iq:m2:how-bonds-trade:2}
Le système de FINRA qui recueille les déclarations de chaque transaction de gré à gré sur les obligations
éligibles, dans les quinze minutes, et les publie. Il affiche le prix, l'heure et la taille, mais les tailles des grandes
transactions sont plafonnées, il n'affiche ni cotations ni profondeur, et il ne montre les transactions
qu'une fois qu'elles ont eu lieu.
\iqlookfor{ce qui est publié et ce qui est masqué.}
\end{solution}

\begin{solution}{iq:m2:how-bonds-trade:3}
Modéliser les cours acheteurs comme une valeur commune diminuée d'une marge, plus des conditions privées indépendantes~; le meilleur
de $n$ s'améliore de $\sigma e_n$~; chaque perdant coûte une fuite $\ell$~; minimiser
$m - \sigma e_n + \ell(n-1)$. La réponse croît avec la dispersion des conditions des dealers
et décroît avec la fuite.
\iqlookfor{le compromis et sa statique comparative.}
\end{solution}

\begin{solution}{iq:m2:how-bonds-trade:4}
Valoriser chaque obligation à partir des \omterm{def:m2:how-bonds-trade:composite}{prix composites} et des courbes des émetteurs~; agréger le risque du panier
(duration, duration de spread, secteurs, liquidité)~; tarifer le coût de sa couverture au moyen de
swaps d'indices et d'ETF, et celui de la liquidation des obligations sur plusieurs jours, obligation par obligation selon
leur liquidité~; ajouter une charge pour le risque résiduel et le capital du panier~; coter un
prix unique.
\iqlookfor{risque au niveau du panier, instruments de couverture et coût de liquidation.}
\end{solution}

\begin{solution}{iq:m2:how-bonds-trade:5}
Sa valeur liquidative repose sur des cours d'obligations pouvant être périmés ou estimés, alors que
l'ETF se négocie en continu~; en période de tension, les vendeurs passent par l'ETF, dont le cours baisse en premier.
L'arbitrage des participants autorisés est limité par le coût et le risque de négociation des
obligations. Le cours de l'ETF est souvent la meilleure mesure du prix auquel les obligations pourraient
réellement se vendre.
\iqlookfor{valeur liquidative périmée et limites de l'arbitrage.}
\end{solution}

\begin{solution}{iq:m2:how-bonds-trade:6}
Précalculer les \omterm{def:m2:how-bonds-trade:composite}{prix composites} et les sensibilités de couverture de chaque obligation, et les mettre à jour à
chaque événement de marché~; à chaque RFQ, retrouver l'obligation, appliquer les ajustements liés au client, à la taille et à l'inventaire
à partir de tables en cache, vérifier les limites, puis répondre~; garder le chemin d'exécution exempt d'E/S
et d'allocation~; journaliser les décisions pour une analyse ultérieure~; suivre les taux de réponse et les mark-outs pour
recalibrer.
\iqlookfor{précalcul, recherches rapides, limites, retour d'information.}
\end{solution}
